% a simple template with cesbio color code

\documentclass[xcolor=table, t, aspectratio=169]{beamer} 
\graphicspath{{../FIGURES/}}
\usepackage[utf8]{inputenc}
\usepackage{booktabs, comment} 
\usepackage[absolute, overlay]{textpos} 
\usepackage{pgfpages}
\usepackage{caption}
\usepackage[font=footnotesize]{caption}
\usepackage{multirow}
\usepackage{hyperref}
\usepackage{blindtext}
\usepackage{amsmath}
\usepackage[makeroom]{cancel}
\usepackage{textpos}
\usepackage{tikz}
\usepackage{csquotes}
\usepackage{graphicx}


% CESBIO ADAPTED COLOURS 
\definecolor{colorLAERO}{RGB}{67, 141, 204}
\definecolor{forestgreen}{rgb}{0.13, 0.55, 0.13}
\definecolor{greenyellow}{rgb}{0.68, 1.0, 0.18}
\definecolor{inchworm}{rgb}{0.7, 0.93, 0.36}

% OPTIONAL COLOURS
% you can add the ones you need : lot of proposition on http://latexcolor.com/
\definecolor{ballblue}{rgb}{0.13, 0.67, 0.8}
\definecolor{bluegreen}{rgb}{0.0, 0.87, 0.87}
\definecolor{chromeyellow}{rgb}{1.0, 0.65, 0.0}
\definecolor{darkcyan}{rgb}{0.0, 0.55, 0.55}
\definecolor{darkpastelred}{rgb}{0.76, 0.23, 0.13}
\definecolor{electricgreen}{rgb}{0.0, 1.0, 0.0}
\definecolor{electriclime}{rgb}{0.8, 1.0, 0.0}
\definecolor{deepcarrotorange}{rgb}{0.91, 0.41, 0.17}
\definecolor{ganttblue}   {HTML}{2B6CB0}
\definecolor{ganttcyan}   {HTML}{2C7A7B}
\definecolor{ganttgreen}  {HTML}{276749}
\definecolor{ganttorange} {HTML}{C05621}
\definecolor{ganttgray}   {HTML}{4A5568}
\definecolor{ganttbg}     {HTML}{EBF4FF}
\definecolor{ganttbar}    {HTML}{BEE3F8}
\definecolor{milestonecol}{HTML}{E53E3E}


% CESBIO ADAPTED BEAMER OPTIONS
\setbeamercolor{title in head/foot}{bg=inchworm, fg=forestgreen}
\setbeamercolor{author in head/foot}{bg=colorLAERO}
\setbeamertemplate{page number in head/foot}{}

\useoutertheme{infolines}
\usetheme{Warsaw}
\usecolortheme[named=colorLAERO]{structure}


% SPECIFIC COMMANDS
% if you nee to set command for youe presentation, you can put them here
\newcommand{\defnotation}[1]{\bf{\textsl{#1}}}

\usepackage[
citestyle=abbrv, 
bibstyle=abbrv,
style=authoryear, %-comp,
backend=biber
%maxbibnames=99, % Make sure we are printing all authors in the appendix
]{biblatex}

\usepackage{pgfgantt}
\usepackage{helvet}
% PRESENTATION INFORMATIONS
% to fill in

% \title[]{}
% \subtitle{Subtitle}

\usetikzlibrary{positioning, calc}
\usetikzlibrary{calc}

% ── Colours (adapt to your theme) ─────────────────────────────────────────────
\definecolor{titleblue}  {HTML}{1A365D}
\definecolor{accentblue} {HTML}{2B6CB0}
\definecolor{lightgray}  {HTML}{EDF2F7}
\definecolor{medgray}    {HTML}{718096}

% ── Metadata ──────────────────────────────────────────────────────────────────
\author[J.~Hedman]{Johan Hedman}
\institute[LAERO]{Laboratoire d'Aérologie (LAERO) · Université de Toulouse}
\date{\today}

% ── Title frame ───────────────────────────────────────────────────────────────
% Use a fully custom frame instead of \titlepage for full layout control
\newcommand{\customtitleframe}{%
\begin{frame}[plain]
	\begin{tikzpicture}[remember picture, overlay]
		
		%% ── Background ─────────────────────────────────────────────────────────────
		\fill[titleblue]
		(current page.north west) rectangle
		($(current page.north east) + (0, -0.38\paperheight)$);
		
		\fill[accentblue]
		($(current page.north west) + (0, -0.38\paperheight)$) rectangle
		($(current page.north east) + (0, -0.405\paperheight)$);
		
		%% ── Title block ────────────────────────────────────────────────────────────
		\node[anchor=west, text=white, font=\large\bfseries, text width=0.95\paperwidth,] at 
		($(current page.north west) + (0.7cm, -0.13\paperheight)$)
		{Etude des rouleaux de couche limite sous les tempêtes};
		
		\node[anchor=west, text=white!70!titleblue, font=\small,text width=0.85\paperwidth,] at 
		($(current page.north west) + (0.7cm, -0.26\paperheight)$)
		{Johan Hedman \quad·\quad LAERO, Université de Toulouse \quad·\quad \insertdate};
		
		%% ── Supervisors ────────────────────────────────────────────────────────────
		\node[anchor=north west, font=\bfseries\small\color{accentblue}] at 
		($(current page.north west) + (0.7cm, -0.405\paperheight - 0.52cm)$)
		{Supervision};
		
		\node[anchor=north west, font=\footnotesize\color{titleblue}, align=left, text width=0.38\paperwidth]
		at ($(current page.north west) + (0.7cm, -0.405\paperheight - 1.0cm)$)
		{%
			\textbf{Prof. Jean-Pierre Chaboureau} \\
			\textcolor{medgray}{LAERO, Université de Toulouse} \\[5pt]
		};
		
		%% ── Divider ────────────────────────────────────────────────────────────────
		\draw[accentblue!30, line width=0.5pt]
		($(current page.north west) + (0.50\paperwidth, -0.405\paperheight - 0.50cm)$) --
		($(current page.south west) + (0.50\paperwidth, 0.45cm)$);
		
		%% ── Jury ───────────────────────────────────────────────────────────────────
		\node[anchor=north west, font=\bfseries\small\color{accentblue}]at ($(current page.north west) + (0.53\paperwidth, -0.405\paperheight - 0.52cm)$)
		{Membre du CSI};
		
		\node[anchor=north west, font=\footnotesize\color{titleblue}, align=left, text width=0.42\paperwidth]
		at ($(current page.north west) + (0.53\paperwidth, -0.405\paperheight - 1.0cm)$)
		{%
			\textbf{Prof } \textcolor{medgray}{\textit{}} \\
			\textbf{Prof.~Firstname Lastname} \textcolor{medgray}{\textit{}} \\
		};
		
		%% ── Logo strip ─────────────────────────────────────────────────────────────
		\fill[white]
		($(current page.south west) + (0, 0.10\paperheight)$)
		rectangle (current page.south east);
		
		\draw[accentblue!25, line width=0.4pt]
		($(current page.south west) + (0, 0.10\paperheight)$) --
		($(current page.south east) + (0, 0.10\paperheight)$);
		
		\node[anchor=south west, xshift=0.5cm, yshift=0.20cm] at (current page.south west)		{
\includegraphics[height=1.5cm]{LOGO/LAERO-logo.jpg}};
		
		\node[anchor=south, yshift=0.25cm] at (current page.south)
			{
\includegraphics[height=1.0cm]{LOGO/UNIV-logo.png}\hspace{1cm}%
			
\includegraphics[height=1.20cm]{LOGO/CNRS-logo.jpeg}};
		
		\node[anchor=south east, xshift=-0.5cm, yshift=0.20cm] at (current page.south east)
		{
\includegraphics[height=1.2cm]{LOGO/logo-IRD.png}};
		
	\end{tikzpicture}
\end{frame}
}





% cesbio adapted frame structure
\addtobeamertemplate{navigation symbols}{}{%
    \usebeamerfont{footline}%
    \usebeamercolor[fg]{footline}%
    \hspace{1em}%
    \insertframenumber/\inserttotalframenumber
}

\usepackage{soul}
\makeatletter
\let\HL\hl
\renewcommand\hl{%
  \let\set@color\beamerorig@set@color
  \let\reset@color\beamerorig@reset@color
  \HL}
\makeatother

%\addbibresource{Rolls.bib}
\graphicspath{{../../../FIGURES/}}

%Title Page
\title{Réunion Jeudi 20/11/2025}
\author{}

\usepackage[backend=biber,style=authoryear]{biblatex}
\addbibresource{../../groupeThese.bib}

\begin{document}
\begin{frame}
\maketitle
\end{frame}

\begin{frame}{Le projet}
\textbf{Mes buts :}

\begin{itemize}
  \item Reproduire des rouleaux dans un \textbf{cadre idéalisé} (travail de Wahiba et Thomas)
  \item Influence des paramètres (profil de température, flux de surface, hétéréogénéité des flux, vents) sur les rouleaux (taille, formation, durée de vie)
  \item Evaluer transports de moments par les rouleaux ? (\cite{gryschkaImpactForcedRoll2014}) 
  \item Paramétrisation des rouleaux ? (\cite{fosterWhyRollsAre2005},  \cite{wangDeepLearningParameterization2023})
  \item Avoir descriptions moodèle forward des rouleaux pour ensuite faire le problème inverse et récupérer des flux ? 
  \end{itemize}
  \textbf{Etat d'avancement:}
  \begin{itemize}
    \item  Test pour obtenir des rouleaux avec des domaines idéeaux et simples. Lecture de littérature. 
  \end{itemize}
\end{frame}

\section{Experience noVAPOR}
\begin{frame}{Simulation noVAPOR}
 \begin{columns}
 \begin{column}{0.48\textwidth}
 \begin{itemize}
	\item Air sec, pas de vapeur d'eau. 
	\item Forte inversion de température - vents proche de Wahiba 
  \item Bords ouverts 
	\item Grand domaine pour avoir formation de rouleaux mais couteux
\end{itemize}
\end{column}
\begin{column}{0.48\textwidth}
\begin{figure}
	\includegraphics[width=0.4\textheight]{noVAPOR/inputs.png}
\end{figure}
\end{column}
\end{columns}
\end{frame}

\subsection{Profiles vitesses}
\begin{frame}{Profiles}
Pb vents verticaux en altitude + CL se stabilise (peut etre du a temperétaure au sol plus faible que profil de temp dans mNH)
	\begin{figure}[t]
		\centering
		\includegraphics[width=0.7\textwidth]{noVAPOR/wind_profiles.png}
	\end{figure}
\end{frame}

\subsection{Flux}
\begin{frame}{Flux}
Flux sensible cohérent avec papier de Wahiba ; un peu trop forts ($\approx 100W.m^{-2}$)
	\begin{figure}[t]
		\centering
		\includegraphics[width=0.7\textwidth]{noVAPOR/fluxes.png}
	\end{figure}
\end{frame}

\subsection{WT}
\begin{frame}{Vitesses verticales}
	petits rouleaux (dû à des flux trop forts ? ressemble à des free rolls \> pas d'influence des flux ?) faire tracé de f ? 
	\begin{figure}[t]
		\centering
		\includegraphics[width=0.7\textwidth]{noVAPOR/WT_ni350_t20.png}
	\end{figure}
\end{frame}


\section{Experience miniROLLS}
\begin{frame}{Simulation miniROLLS}
 \begin{columns}
 \begin{column}{0.48\textwidth}
 \begin{itemize}
	\item remise bords cyclique pour accélerer le CT (conditionne surement le nombre de rouleaux mais a étudier quand on en a)
	\item mise de profils exacts de Wahiba (par contre on enleve la forte inversion )
	\item $\Delta x = \Delta y = 100m$ ; $\Delta t = 0.5 s$ pou temps de calcul correct
\end{itemize}
\end{column}
\begin{column}{0.48\textwidth}
\begin{figure}
	\includegraphics[width=0.4\textheight]{miniROLLS/inputs.png}
\end{figure}
\end{column}
\end{columns}
\end{frame}

\subsection{Profiles vitesses}
\begin{frame}{Profiles}
Pb vents verticaux en altitude + CL se stabilise (peut etre du a temperétaure au sol plus faible que profil de temp dans mNH)
	\begin{figure}[t]
		\centering
		\includegraphics[width=0.7\textwidth]{miniROLLS/wind_profiles.png}
	\end{figure}
\end{frame}

\subsection{Flux}
\begin{frame}{Flux}
Flux plus cohérent avec Wahiba et stables; tjs couche qui se stabilise avec un sol plus froid que la temperature du niveau 0 \hl{a changer}
	\begin{figure}[t]
		\centering
		\includegraphics[width=0.7\textwidth]{miniROLLS/fluxes.png}
	\end{figure}
\end{frame}

\subsection{WT}
\begin{frame}{Vitesses verticales}
  Vitesses verticales plus faible mais bruit : problème au bord droit qui se propage toujours ? 
	\begin{figure}[t]
		\centering
		\includegraphics[width=0.7\textwidth]{miniROLLS/WT_ni100_t20.png}
	\end{figure}
\end{frame}
\end{frame}
\begin{frame}{}
\printbibliography

\end{frame}


\end{document}   
